%% ma: L10-B14-0091
%% lop: 10
%% chuong: 5
%% bai: 14
%% dang: 10.14.12
%% loai: DS
%% muc: [NB, NB, TH, TH]
%% dap_an: "SSĐS"
%% nguon: "chuyen-de-cac-so-dac-trung-cua-mau-so-lieu-khong-ghep-nhom-toan-10-chuong-trinh-moi.pdf (GV. Phan Nhật Linh) (P05), Bài 03 Dạng 2 Phần II Câu 8"
%% kien_thuc: "Lớp 10 Bài 13 (số trung bình, trung vị, tứ phân vị, mốt), Bài 14 (khoảng biến thiên, khoảng tứ phân vị, phương sai, độ lệch chuẩn, giá trị bất thường)"
%% trang_thai: chinh-thuc
%% ghi_chu: "Tự duyệt theo yêu cầu giáo viên 22/10/2026."
\begin{ex}
Sản lượng lúa (đơn vị: tạ) của $40$ thửa ruộng thí nghiệm có cùng diện tích được trình bày trong bảng tần số sau:
\begin{center}
\begin{tabular}{|c|c|c|c|c|c|c|}\hline
Sản lượng $(x)$ & $20$ & $21$ & $22$ & $23$ & $24$ & \\\hline
Tần số $(n)$ & $5$ & $8$ & $11$ & $10$ & $6$ & $n=40$\\\hline
\end{tabular}
\end{center}
Xét tính đúng sai của các khẳng định sau:
\choiceTF{Mốt của mẫu số liệu là $11$}{Số trung bình của mẫu số liệu là $\overline x=23$}{\True Trung vị của mẫu số liệu là $M_e=22$}{Phương sai của mẫu số liệu là $s^2=2{,}04$}
\loigiai{a) Sai: tần số lớn nhất là $11$ ứng với giá trị $22$, nên $M_o=22$.

b) Sai: $\overline x=22{,}1$.

c) Đúng: $M_e=\dfrac{x_{20}+x_{21}}2=22$.

d) Sai: $s^2=\dfrac{61{,}6}{40}=1{,}54$.}
\end{ex}
